\documentclass[10pt,a4paper]{amsart}
\usepackage{amsmath,amssymb,amsthm,graphicx}
% Retained minimal preamble for the Stacks Project chapter more-algebra.tex (More on Algebra).
% The chapter file pinned-original-source.tex carries only \input{preamble} on its first line, so no
% preamble text lives inside the chapter itself.  The Stacks Project uses one single shared preamble for
% every chapter; the snapshot stages it identically as sources/stacks-chow/preamble.tex and
% sources/stacks-etale-cohomology/preamble.tex.  The block below is copied mechanically, line for line,
% from sources/stacks-chow/preamble.tex lines 14-18 (the Xy-pic setup) and 154-217 (the native theorem
% declarations together with the mathematical macros).  Nothing was added, removed or redefined; the
% wrapper class and the amsmath/amssymb/amsthm/graphicx base packages are supplied by the assembler.
% Chapter: more-algebra.tex -- More on Algebra.
% For commutative diagrams we use Xy-pic
\usepackage[all]{xy}

% We use 2cell for 2-commutative diagrams.
\xyoption{2cell}

% Theorem environments.
%
\theoremstyle{plain}
\newtheorem{theorem}[subsection]{Theorem}
\newtheorem{proposition}[subsection]{Proposition}
\newtheorem{lemma}[subsection]{Lemma}

\theoremstyle{definition}
\newtheorem{definition}[subsection]{Definition}
\newtheorem{example}[subsection]{Example}
\newtheorem{exercise}[subsection]{Exercise}
\newtheorem{situation}[subsection]{Situation}

\theoremstyle{remark}
\newtheorem{remark}[subsection]{Remark}
\newtheorem{remarks}[subsection]{Remarks}

\numberwithin{equation}{subsection}

% Macros
%
\def\lim{\mathop{\mathrm{lim}}\nolimits}
\def\colim{\mathop{\mathrm{colim}}\nolimits}
\def\Spec{\mathop{\mathrm{Spec}}}
\def\Hom{\mathop{\mathrm{Hom}}\nolimits}
\def\Ext{\mathop{\mathrm{Ext}}\nolimits}
\def\SheafHom{\mathop{\mathcal{H}\!\mathit{om}}\nolimits}
\def\SheafExt{\mathop{\mathcal{E}\!\mathit{xt}}\nolimits}
\def\Sch{\mathit{Sch}}
\def\Mor{\mathop{\mathrm{Mor}}\nolimits}
\def\Ob{\mathop{\mathrm{Ob}}\nolimits}
\def\Sh{\mathop{\mathit{Sh}}\nolimits}
\def\NL{\mathop{N\!L}\nolimits}
\def\CH{\mathop{\mathrm{CH}}\nolimits}
\def\proetale{{pro\text{-}\acute{e}tale}}
\def\etale{{\acute{e}tale}}
\def\QCoh{\mathit{QCoh}}
\def\Ker{\mathop{\mathrm{Ker}}}
\def\Im{\mathop{\mathrm{Im}}}
\def\Coker{\mathop{\mathrm{Coker}}}
\def\Coim{\mathop{\mathrm{Coim}}}

% Boxtimes
%
\DeclareMathSymbol{\boxtimes}{\mathbin}{AMSa}{"02}

%
% Macros for moduli stacks/spaces
%
\def\QCohstack{\mathcal{QC}\!\mathit{oh}}
\def\Cohstack{\mathcal{C}\!\mathit{oh}}
\def\Spacesstack{\mathcal{S}\!\mathit{paces}}
\def\Quotfunctor{\mathrm{Quot}}
\def\Hilbfunctor{\mathrm{Hilb}}
\def\Curvesstack{\mathcal{C}\!\mathit{urves}}
\def\Polarizedstack{\mathcal{P}\!\mathit{olarized}}
\def\Complexesstack{\mathcal{C}\!\mathit{omplexes}}
% \Pic is the operator that assigns to X its picard group, usage \Pic(X)
% \Picardstack_{X/B} denotes the Picard stack of X over B
% \Picardfunctor_{X/B} denotes the Picard functor of X over B
\def\Pic{\mathop{\mathrm{Pic}}\nolimits}
\def\Picardstack{\mathcal{P}\!\mathit{ic}}
\def\Picardfunctor{\mathrm{Pic}}
\def\Deformationcategory{\mathcal{D}\!\mathit{ef}}

% --- appended 2026-10-12 (Chief): the project own definitions of the reference/slogan/history
% display environments, copied mechanically from the shared Stacks preamble lines 8-12
% (\usepackage{verbatim} plus the three \newenvironment lines, which the chapter preamble keeps
% verbatim from the project preamble). Nothing added, removed or redefined: these environments
% carry no mathematical content and the project itself typesets them as comments. Without these
% three lines a counted statement whose own text uses them fails with Environment slogan undefined.
\usepackage{verbatim}
\newenvironment{reference}{\comment}{\endcomment}
\newenvironment{slogan}{\comment}{\endcomment}
\newenvironment{history}{\comment}{\endcomment}
\usepackage{cleveref}
\crefname{originallemma}{original lemma}{original lemmas}
\Crefname{originallemma}{Original lemma}{Original lemmas}
\makeatletter
\newlabel{lemma-K-flat-resolution}{{06Y4}{source}{}{sourcepointer-1}{}}
\newlabel{lemma-K-flat-resolution@cref}{{[originallemma][0][]06Y4}{[1][0][]source}{}{}{}}
\makeatother
\crefname{originallemma}{original lemma}{original lemmas}
\Crefname{originallemma}{Original lemma}{Original lemmas}
\makeatletter
\newlabel{lemma-base-change-K-flat}{{06Y1}{source}{}{sourcepointer-2}{}}
\newlabel{lemma-base-change-K-flat@cref}{{[originallemma][0][]06Y1}{[1][0][]source}{}{}{}}
\makeatother
\crefname{originallemma}{original lemma}{original lemmas}
\Crefname{originallemma}{Original lemma}{Original lemmas}
\makeatletter
\newlabel{algebra-lemma-flat-base-change}{{00HI}{source}{}{sourcepointer-3}{}}
\newlabel{algebra-lemma-flat-base-change@cref}{{[originallemma][0][]00HI}{[1][0][]source}{}{}{}}
\makeatother
\crefname{originallemma}{original lemma}{original lemmas}
\Crefname{originallemma}{Original lemma}{Original lemmas}
\makeatletter
\newlabel{lemma-tensor-product-K-flat}{{0795}{source}{}{sourcepointer-4}{}}
\newlabel{lemma-tensor-product-K-flat@cref}{{[originallemma][0][]0795}{[1][0][]source}{}{}{}}
\makeatother
\crefname{originallemma}{original lemma}{original lemmas}
\Crefname{originallemma}{Original lemma}{Original lemmas}
\makeatletter
\newlabel{lemma-K-flat-two-out-of-three}{{06Y2}{source}{}{sourcepointer-5}{}}
\newlabel{lemma-K-flat-two-out-of-three@cref}{{[originallemma][0][]06Y2}{[1][0][]source}{}{}{}}
\makeatother
\title{Base change along the map sending a variable to zero splits the complex off its own shift}
\author{The Stacks Project Authors}
\date{stacks-project:more-algebra.tex:local-snapshot-20261012; GNU FDL 1.2 or later}
\begin{document}
\maketitle
\setcounter{subsection}{1}
\noindent\textit{Editorial note.} All delivered mathematical text is the author's own, copied from the pinned Stacks Project chapter file more-algebra.tex; no statement or proof step was rewritten, reordered or removed. The counted statement is the labeled lemma environment at source lines 21217-21226 (canonical tag 065F) and the counted proof is the proof environment at source lines 21228-21256, which immediately follows it and is its own proof.\par\medskip
\noindent\textit{Editorial note.} The proof displays the final computation as a chain of equalities ending with the direct sum of the resolution and its shift; it is copied exactly as printed, including the use of the total complex notation of the shared preamble.\par\medskip
\noindent\textit{Editorial note.} The five references are four earlier results of this chapter on K-flat resolutions and four of the five citations concern flatness and K-flatness, plus a base change statement of another chapter; all five are navigation stubs with their roles recorded in the reasons.\par\medskip
\noindent\textit{Editorial note.} The preamble span runs to line 90, the current end of the retained preamble, which includes the appended environment definitions; the counted text does not use them, and the span is taken whole so that the file version is unambiguous.\par\medskip
\noindent\textit{Editorial note.} The source URL is the source-level Stacks Project URL, as required when more than one card is produced from one source; the per-claim tag is carried by claim\_key and by this note. The counted statement and proof cite no bibliography entry, so this card needs no bibliography.\par\medskip
\noindent\textit{Allowed standard prerequisites (editorial).}\begin{itemize}
\item Every complex of modules admits a resolution by a K-flat complex whose terms are flat modules. invoked at proof lines 21228-21230. Provides the representative of the complex on which the computation is performed.
\item Extending scalars along a flat algebra change preserves K-flatness and flatness of the terms, and the total complex of a double complex built from a K-flat resolution is again K-flat. used at proof lines 21230-21246. Shows that the total complex used in the display is a K-flat complex over the polynomial algebra, so that it can be used to compute the derived tensor product.
\item Multiplication by the variable on the extended complex is injective with cokernel isomorphic to the original complex, so the associated total complex is quasi-isomorphic to the original complex. used at proof lines 21235-21240. Identifies the total complex as a resolution of the original complex, which justifies replacing the complex by it.
\end{itemize}\medskip
\ifdefined\hypertarget\hypertarget{sourcepointer-1}{}\fi
\noindent\textit{Source reference.} Original lemma 06Y4: sources/stacks-more-algebra/pinned-original-source.tex, source lines 14380--14380. The reference supplies the K-flat resolution by flat modules chosen at proof lines 21228-21230; the double complex construction and the computation are inside the counted span.\par\medskip
\ifdefined\hypertarget\hypertarget{sourcepointer-2}{}\fi
\noindent\textit{Source reference.} Original lemma 06Y1: sources/stacks-more-algebra/pinned-original-source.tex, source lines 14210--14210. The reference is invoked at proof lines 21231-21233 to see that the extended complex is K-flat with flat terms; the total complex argument built on it is inside the counted span.\par\medskip
\ifdefined\hypertarget\hypertarget{sourcepointer-3}{}\fi
\noindent\textit{Source reference.} Original lemma 00HI: sources/stacks-chow/tags.txt, source lines 644--644. The reference is cited together with the previous one at proof lines 21232-21233 for the flatness of the extended terms; it is a standard pointer and its use is inside the counted span.\par\medskip
\ifdefined\hypertarget\hypertarget{sourcepointer-4}{}\fi
\noindent\textit{Source reference.} Original lemma 0795: sources/stacks-more-algebra/pinned-original-source.tex, source lines 14224--14224. The reference is one of the two cited at proof lines 21244-21246 to conclude that the associated total complex is K-flat; the computation that uses that K-flatness is inside the counted span.\par\medskip
\ifdefined\hypertarget\hypertarget{sourcepointer-5}{}\fi
\noindent\textit{Source reference.} Original lemma 06Y2: sources/stacks-more-algebra/pinned-original-source.tex, source lines 14241--14241. The reference is the alternative second citation at proof lines 21244-21246 for the same conclusion; it is a standard pointer offered as an option by the author and its use is inside the counted span.\par\medskip
\begin{lemma}
\label{lemma-pull-push}
Let $R$ be a ring. Let $K^\bullet$ be a complex of $R$-modules.
Consider the $R$-algebra map $R[x] \to R$ which maps $x$ to zero.
Then
$$
K^\bullet \otimes_{R[x]}^{\mathbf{L}} R \cong K^\bullet \oplus K^\bullet[1]
$$
in $D(R)$.
\end{lemma}

\begin{proof}
Choose a K-flat resolution $P^\bullet \to K^\bullet$ over $R$
such that $P^n$ is a flat $R$-module for all $n$, see
Lemma \ref{lemma-K-flat-resolution}. Then $P^\bullet \otimes_R R[x]$
is a K-flat complex of $R[x]$-modules whose terms are flat $R[x]$-modules, see
Lemma \ref{lemma-base-change-K-flat} and
Algebra, Lemma \ref{algebra-lemma-flat-base-change}.
In particular $x : P^n \otimes_R R[x] \to P^n \otimes_R R[x]$
is injective with cokernel isomorphic to $P^n$. Thus
$$
P^\bullet \otimes_R R[x] \xrightarrow{x} P^\bullet \otimes_R R[x]
$$
is a double complex of $R[x]$-modules whose associated
total complex is quasi-isomorphic to $P^\bullet$ and hence $K^\bullet$.
Moreover, this associated total complex is a K-flat complex
of $R[x]$-modules for example by
Lemma \ref{lemma-tensor-product-K-flat} or by
Lemma \ref{lemma-K-flat-two-out-of-three}.
Hence
\begin{align*}
K^\bullet \otimes_{R[x]}^{\mathbf{L}} R
& \cong
\text{Tot}(P^\bullet \otimes_R R[x] \xrightarrow{x} P^\bullet \otimes_R R[x])
\otimes_{R[x]} R =
\text{Tot}(P^\bullet \xrightarrow{0} P^\bullet) \\
& = P^\bullet \oplus P^\bullet[1] \cong K^\bullet \oplus K^\bullet[1]
\end{align*}
as desired.
\end{proof}

\end{document}
